\section{Pacientes Vinculados y Revocaciones}

Tal como se enfatizó previamente, el portal web del equipo médico carece deliberadamente de una infraestructura de agregación que permita invitar por nombres al paciente. El software confiere \textbf{soberanía rotunda de los permisos clínicos al paciente original}, y son ellos quienes dictan la dinámica en su plataforma, mientras esto produce un impacto asimétrico y directo en el tuyo.
Esta arquitectura ha sido minuciosamente pensada cumpliendo los mejores estándares de cuidado y derecho humano de la salud.

\subsection{Recepción y Notificaciones}

Cuando logras concretar una consulta presencial con un cliente/paciente que cuenta con estudios en nuestra sucursal y este en su portal, dentro de su sección pertinente realiza una "Búsqueda" exitosa de tus apellidos/cédula en nuestra red activa aprobada y oprime "\textbf{Agregar o Vincular Médico}":
\begin{enumerate}
    \item Automáticamente y sin que realices ninguna acción de gestión en tu pantalla, los registros centralizados te crearán una "alianza de lectura".
    \item Verás su nombre rebotar o aparecer en el instante de refresco, agregando un elemento numérico nuevo a tu bloque total de pacientes desde el Dashboard de "Pacientes Vinculados".
    \item A partir de ese momento, contarás con una vía hermenéutica compartida para todos los historiales pasados (estudios antiguos almacenados) e historiales futuros, siempre que la vinculación permanezca activa.
\end{enumerate}

\subsection{Entendiendo las Revocaciones}

¿Qué significa que el paciente es soberano sobre su vinculación? Qué en cualquier momento, puede retirar abruptamente el consentimiento, y debes estar prevenido para esa eventualidad lógica:
\begin{enumerate}
    \item Si tu paciente decide interrumpir su seguimiento de tu lado u otorgar revisión clínica a un distinto doctor alternativo, pueden acceder y oprimir "Revocar o Eliminar Médico".
    \item El software operará un "Borrado Funcional" hacia ti. En cuanto inicies sesión por la mañana (o refresques una pantalla o cambies de ventana), \textbf{dicho paciente y toda la biblioteca subyacente de imágenes/PDFs desaparecen automáticamente y sin dejar rastro informático para tu vista}.
    \item El paciente o los documentos nunca fueron "elminados" o "extraviados erróneamente por nuestro sistema web radiográfico", por lo que si observas una tarjeta perdida en tu base, no levantes un boleto de soporte, ya que es un indicio legítimo de desvinculación a tu perfil por parte de él. 
\end{enumerate}
